\documentclass[12pt,a4paper]{article}
\usepackage[T2A]{fontenc}
\usepackage[utf8]{inputenc}
\usepackage[english, russian]{babel}
\usepackage{indentfirst}
\usepackage{misccorr}
\usepackage{graphicx}
\usepackage{amsmath}
\usepackage{amssymb}
\usepackage{circuitikz}
\usepackage[font={small}]{caption}
\usepackage[left=20mm, top=20mm, right=20mm, bottom=20mm, nohead]{geometry}
\usepackage{float}
\usepackage{tabularx}
\usepackage{array}
\usepackage{longtable}
\usepackage{pstool}
\usepackage{pgfplots}
\usepackage{hhline}
\usepackage{multirow}
\usepackage{wrapfig}
\usepackage{pdfpages}

\DeclareCaptionLabelSeparator{fill}{.\\}

\DeclareCaptionLabelFormat{fullparents}{\bothIfFirst{#1}{~}#2}

\pgfplotsset{compat=1.17}

\begin{document}

    \begin{titlepage}
        \begin{center}
            {\LARGE Отчёт по лабораторной работе 4.7.3.\\}
            \vspace*{11cm}
                \textbf{\LARGE Изучение поляризованного света.}
            \vspace*{6.5cm}
        \end{center}
        \hfill\begin{minipage}{0.37\textwidth}
            Работу выполнил Громов Артём
            \\
            ЛФИ Б02-006
        \end{minipage}
        \vspace{4.8cm}
        \begin{center}
            Долгопрудный, 2022 г.
        \end{center}
    \end{titlepage}

    \section{Аннотация}

        \begin{flushleft}
            \textbf{Цель работы:} ознакомление с методами получения и анализа поляризованного света.
        \end{flushleft}
        \begin{flushleft}
            \textbf{В работе используется:} оптическая скамья с осветителем; зеленый светофильтр; два поляроида; черное зеркало; полированная эбонитовая пластинка;
            стопа стеклянных пластинок; слюдяные пластинки разной толщины; пластинки в 1/4 и 1/2 длины волны; пластинка в одну длину волны для зеленого света
            (пластинка чувствительного оттенка).
        \end{flushleft}

        \par\textbf{Естественный и поляризованный свет.} Как известно, световые волны поперечны: электрический вектор $\mathbf{E}$ и магнитный вектор $\mathbf{H}$
        (или $\mathbf{B}$) взаимно перпендикулярны и располагаются в плоскости, перпендикулярной направлению распространения волны (лучу $\mathbf{S}$). Во всякой
        данной точке пространства ориентация пары векторов $\mathbf{E}$ и $\mathbf{H}$ в плоскости, перпендикулярной лучу $\mathbf{S}$, может, вообще говоря,
        изменяться со временем. В зависимости от характера такого изменения различают естественный и поляризованный свет.
        \par Обычные источники света являются совокупностью огромного числа быстро высвечивающихся ($10^{-10}\div10^{-9}$ с) элементарных источников (атомов или
        молекул), испускающих свет независимо друг от друга, с разными фазами и с разными ориентациями векторов $\mathbf{E}$ и $\mathbf{H}$. Ориентация векторов
        $\mathbf{E}$ и $\mathbf{H}$ в результирующей волне поэтому хаотически изменяется во времени, так что в плоскости, перпендикулярной лучу $\mathbf{S}$, все
        направления оказываются в среднем равноправными. Такой свет называют естественным или неполяризованным.
        \par При помощи специальных приспособлений (поляризаторов), о которых речь будет идти дальше, естественный свет может быть превращен в линейно поляризованный
        (или, как иногда говорят, в плоскополяризованный). В линейно поляризованной световой волне пара векторов $\mathbf{E}$ и $\mathbf{H}$ не изменяет с течением
        времени своей ориентации. Плоскость $\mathbf{E}$, $\mathbf{S}$ называется в этом случае плоскостью колебаний.
        \par Наиболее общим типом поляризации является эллиптическая поляризация. В эллиптически поляризованной световой волне конец вектора $\mathbf{E}$ (в данной
        точке пространства) описывает некоторый эллипс. Линейно поляризованный свет можно рассматривать как частный случай эллиптически поляризованного света, когда
        эллипс поляризации вырождается в отрезок прямой линии; другим частным случаем является круговая поляризация (эллипс поляризации является окружностью).

        \begin{wrapfigure}{r}{0.3\textwidth}
            \centering
            \includegraphics[width=0.3\textwidth]{images/pic1.pdf}
            \caption{Представление световой волны в виде двух линейно поляризованных волн.}
        \end{wrapfigure}
        \par При теоретическом рассмотрении различных типов поляризации часто бывает удобно проектировать вектор $\mathbf{E}$ в некоторой точке пространства на два
        взаимно перпендикулярных направления (рис. 1). В том случае, когда исходная волна была поляризованной, $E_x$ и $E_y$ когерентны между собой и могут быть
        записаны в виде
        \begin{equation}
            \begin{aligned}
                E_x &= {E_x}_0\cos{kz-\omega t}, \\
                E_y &= {E_y}_0\cos{kz-\omega t-\phi},
            \end{aligned}
        \end{equation}
        где амплитуды ${E_x}_0$, ${E_y}_0$, волновой вектор $k$, частота $\omega$ и сдвиг фаз $\phi$ не зависят от времени. Формулы (1) описывают монохроматический
        свет. Немонохроматический свет может быть представлен суммой выражений типа (1) с различными значениями частоты $\omega$.
        \par Ориентация эллипса поляризации определяется отношением амплитуд ${E_y}_0/{E_x}_0$ и разностью фаз $\phi$. В частности, при $\phi=0$, $\pm\pi$ эллипс
        вырождается в отрезок прямой (линейная поляризация). При $\phi=\pm\pi/2$ главные оси эллипса совпадают с осями $x$, $y$. Если при этом отношение амплитуд
        ${E_y}_0/{E_x}_0=1$, эллипс поляризации вырождается в окружность.
        \par В плоскости $z=z_0$ вектор $\mathbf{E}$ волны (1) вращается против часовой стрелки (при наблюдении навстречу волне), если $0<\phi<\pi$. В этом случае
        говорят о левой эллиптической поляризации волны. Если же $\pi<\phi<2\pi$, вращение вектора $\mathbf{E}$ происходит по часовой стрелке, и волна имеет правую
        эллиптическую поляризацию.
        \par В фиксированный момент времени $t=t_0$ концы вектора $\mathbf{E}$ при различных $z$ лежат на винтовой линии. При этом для левой эллиптической
        поляризации образуется левый винт, а для правой --- правый винт.
        \par Неполяризованный свет также может быть разложен на две линейно поляризованные компоненты; однако в этом случае разность фаз $\phi$ испытывает быстрые
        хаотические изменения, так что колебания $E_x$ и $E_y$ оказываются некогерентными.
    
    \section{Экспериментальная установка}

        \par\textbf{Методы получения линейно поляризованного света.} Для получения линейно поляризованного света применяются специальные оптические приспособления
        —-- поляризаторы. Направление колебаний электрического вектора в волне, прошедшей через поляризатор, называется разрешенным направлением поляризатора.
        \par Всякий поляризатор может быть использован для исследования поляризованного света, т. е. в качестве анализатора. Интенсивность $I$ линейно
        поляризованного света после прохождения через анализатор зависит от угла, образованного плоскостью колебаний с разрешенным направлением анализатора:
        \begin{equation}
            I=I_0\cos^2{\alpha}.
        \end{equation}
        Соотношение (2) носит название закона Малюса.
        \par Опишем несколько способов получения плоскополяризованного света.
        \par Отражение света от диэлектрической пластинки. Отраженный от диэлектрика свет всегда частично поляризован. Степень поляризации света, отраженного от
        диэлектрической пластинки в воздух, зависит от показателя преломления диэлектрика $n$ и от угла падения $i$. Как следует из формул Френеля, полная
        поляризация отраженного света достигается при падении под углом Брюстера, который определяется соотношением
        \begin{equation}
            \tg{i}=n.
        \end{equation}
        В этом случае плоскость колебаний электрического вектора в отражен- ном свете перпендикулярна плоскости падения.
        \par Преломление света в стеклянной пластинке. Поскольку отраженный от диэлектрической пластинки свет оказывается частично (или даже полностью)
        поляризованным, проходящий свет также частично поляризуется. Преимущественное направление колебаний электрического вектора в прошедшем свете совпадает с
        плоскостью преломления луча. Максимальная поляризация проходящего света достигается при падении под углом Брюстера. Для увеличения степени поляризации
        преломлённого света используют стопу стеклянных пластинок, расположенных под углом Брюстера к падающему свету.
        \par Преломление света в двоякопреломляющих кристаллах. Некоторые кристаллы обладают свойством двойного лучепреломления. Это связано с различием
        поляризуемости молекул в разных направлениях (диэлектрическая проницаемость $\epsilon$ определяет показатель преломления среды $n$).
        \par Двоякопреломляющий кристалл называют одноосным, если в нём существует одно направление с экстремальным значением $\epsilon$, а в други
        (перпендикулярных) направлениях значения $\epsilon$ одинаковы (тензор диэлектрической проницаемости образует эллипсоид вращения). Направления вдоль осей
        эллипсоида называют главными, одно из них --- c экстремальным значением $\epsilon$ --- оптической осью. Линейно поляризованная волна, в которой вектор
        $\mathbf{E}$ перпендикулярен оптической оси, называется обыкновенной; если же вектор $\mathbf{E}$ имеет проекцию на оптическую ось, это необыкновенная
        волна. Показатели преломления этих волн обозначают через $n_o$ (ординарная волна) и $n_e$ (экстраординарная).
        \par Преломляясь в таких кристаллах, световой луч разделяется на два луча со взаимно перпендикулярными плоскостями колебаний. Отклоняя один из лучей в
        сторону, можно получить плоскополяризованный свет, --- так устроены поляризационные призмы (Николя, Глана).
        \par Поглощение света в дихроических пластинках. У некоторых двоякопреломляющих кристаллов (например, турмалина) коэффициенты поглощения света для двух
        взаимно перпендикулярно поляризованных лучей отличаются настолько сильно, что уже при небольшой толщине кристалла одни из лучей гасятся практически
        полностью, и из кристалла выходит линейно поляризованный пучок света. Это явление носит название дихроизма. В настоящее время дихроические пластинки
        изготовляются в виде тонких пленок, закрытых стёклами --- поляроидов.
        \par Если поляроиды используются для создания поляризованного света, их называют поляризаторами, если для определения характера поляризации —--
        анализаторами.
        \par\textbf{Определение направления разрешенной плоскости колебаний поляроида.} Определить направление разрешенных колебаний поляроида проще всего с
        помощью черного зеркала.
        \par При падении света на отражающую поверхность под углом Брюстера свет в отражённом луче полностью поляризован, а вектор $\mathbf{E}$ параллелен
        отражающей поверхности («правило иголки»). Луч света, прошедший поляроид и отразившийся от чёрного зеркала, имеет минимальную интенсивность при выполнении
        двух условий: во-первых, свет падает на отражающую поверхность под углом Брюстера, и во-вторых, в падающем пучке вектор $\mathbf{E}$ лежит в плоскости
        падения.
        \par Вращая поляроид вокруг направления луча и чёрное зеркало вокруг оси, перпендикулярной лучу, методом последовательных приближений можно добиться
        минимальной яркости луча, отражённого от зеркала, и таким образом определить разрешённое направление поляроида.
        \par Измеряя угол поворота зеркала (угол Брюстера), нетрудно определить коэффициент преломления материала, из которого изготовлено зеркало. Описанный
        метод часто используется для измерения коэффициента преломления непрозрачных диэлектриков.

        \begin{wrapfigure}{r}{0.3\textwidth}
            \centering
            \includegraphics[width=0.3\textwidth]{images/pic2.pdf}
            \caption{Разложение линейно поляризованного света по главным направлениям двоякопреломляющей пластинки.}
        \end{wrapfigure}
        \par\textbf{Получение эллиптически поляризованного света.} Эллиптически поляризованный свет можно получить из линейно поляризованного с помощью
        двоякопреломляющих кристаллических пластинок.
        \par Двоякопреломляющая пластинка имеет два взаимно перпендикулярных главных направления, совпадающих с осями эллипсоида диэлектрической проницаемости.
        Волны, поляризованные вдоль главных направлений, распространяются в пластинке с разными скоростями, не изменяя характера своей поляризации. Эти волны
        называются главными. Мы будем обозначать показатели преломления для главных волн через $n_{\xi}$ и $n_{\eta}$, где $\xi$ и $\eta$ --- главные направления
        кристаллической пластинки (рис. 2).
        \par Пусть на пластинку падает линейно поляризованная волна, электрический вектор которой ориентирован под некоторым углом $\alpha$ к оси $\xi$. Разложим
        вектор $\mathbf{E}$ на составляющие $E_{\xi}$ и $E_{\eta}$. На входе пластинки $E_{\xi}$ и $E_{\eta}$ находятся в фазе. На выходе из-за разности скоростей
        между ними появляется разность хода (выразим её в долях длины волны, например: $\lambda/2$, $\lambda/3$, ...)
        $$
            \frac{\lambda}{m}=d(n_{\xi}-n_{\eta}),
        $$
        при этом сдвиг фаз определяется соотношением
        \begin{equation}
            \Delta\phi=\frac{2\pi}{m}=kd(n_{\xi}-n_{\eta}),
        \end{equation}
        где $k$ --- волновое число для вакуума, $d$ --— толщина кристаллической пластинки. Как уже отмечалось, при сложении двух взаимно перпендикулярных колебаний,
        обладающих некоторым сдвигом фаз, образуется колебание, поляризованное по эллипсу.
        \par Рассмотрим практически важные частные случаи.
        \par а) Пластинка дает сдвиг фаз $2\pi$ (пластинка в длину волны). В результате сложения волн на выходе пластинки образуется линейно поляризованная волна с
        тем же направлением колебаний, что и в падающей волне.
        
        \begin{wrapfigure}{r}{0.3\textwidth}
            \centering
            \includegraphics[width=0.3\textwidth]{images/pic3.pdf}
            \caption{Поворот направления колебаний с помощью пластинки в $\lambda/2$.}
        \end{wrapfigure}
        \par б) Пластинка дает сдвиг фаз $\pi$ (пластинка в полдлины волны). На выходе пластинки снова образуется линейно поляризованная волна. Направление $bb'$
        колебаний этой волны повернуто относительно направления $aa'$ колебаний падающей волны (рис. 3). Как нетрудно сообразить, направление $bb'$ является
        зеркальным отображением направления $aa'$ относительно одного из главных направлений пластинки. Такую пластинку используют для поворота направления
        колебаний линейно поляризованного света.
        \par в) Пластинка создает между колебаниями сдвиг фаз $\pi/2$ (пластинка в четверть длины волны). При сложении двух взаимно перпендикулярных колебаний,
        имеющих разность фаз $\pi/2$, образуется эллипс, главные оси которого совпадают с координатными осями $\xi$ и $\eta$. При равенстве амплитуд
        ${E_{\xi}}_0 = {E_{\eta}}_0$ возникает круговая поляризация.
        \par\textbf{Анализ эллиптически поляризованного света.} Анализ эллиптически поляризованного света сводится к нахождению главных осей эллипса поляризации и
        к определению направления вращения электрического вектора.
        \par Главные оси эллипса поляризации определяются с помощью анализатора по максимуму и минимуму интенсивности проходящего света. Направление вращения
        электрического вектора может быть найдено с помощью пластинки в четверть длины волны, для которой известно, какая из главных волн, $E_{\xi}$ или $E_{\eta}$,
        имеет большую скорость распространения (и, соответственно, меньшее значение показателя преломления).
        \par Выберем для определенности координатные оси $\xi$ и $\eta$ на пластинке так, чтобы $n_{\xi}<n_{\eta}$. В этом случае главная волна $E_{\xi}$ имеет
        большую скорость распространения. Поместим такую пластинку на пути эллиптически поляризованного света и совместим главные оси эллипса поляризации с главными
        направлениями пластинки $\lambda/4$. На выходе из этой пластинки сдвиг фаз между $E_{\xi}$ и $E_{\eta}$ вместо $\pi/2$ станет равным нулю или $\pi$. Свет
        окажется линейно поляризованным. Из двух возможных значений сдвига фаз, $0$ или $\pi$, реализуется одно: то, которое соответствует имеющемуся в волне
        направлению вращения электрического вектора.
        \par Рассмотрим, например, случай, когда электрический вектор в эллиптически поляризованной волне вращается против часовой стрелки, если смотреть навстречу
        лучу. В этом случае, очевидно, в волне, падающей на пластинку в $\lambda/4$, колебание $E_{\eta}$ отстает по фазе на $\pi/2$ от колебания $E_{\xi}$. При\
        прохождении через пластинку разность фаз увеличивается до $\pi$. На выходе из пластинки, таким образом, возникают линейно поляризованные волны со сдвигом
        фаз $\pi$. Сложение этих волн дает плоскополяризованную волну, электрический вектор которой располагается во втором и четвертом квадрантах координатной
        системы $\xi$, $\eta$.
        \par Рассуждая аналогичным образом, найдем, что при вращении электрического вектора по часовой стрелке направление колебаний в линейно поляризованной волне,
        выходящей из пластинки, располагается в первом и третьем квадрантах. Определяя направление колебаний на выходе из пластинки с помощью поляроида, можно,
        таким образом, определить характер эллиптической поляризации (вращение против или по часовой стрелке).
        \par\textbf{Пластинка чувствительного оттенка.} Выше предполагалось известным, какому из двух главных направлений пластинки в четверть длины волны
        соответствует большая скорость распространения света. Установить это можно различными способами, например с помощью пластинки чувствительного оттенка (так
        называют пластинку в $\lambda$ для зеленой спектральной компоненты, $\lambda=560$ нм).

        \begin{wrapfigure}{r}{0.3\textwidth}
            \centering
            \includegraphics[width=0.3\textwidth]{images/pic4.pdf}
            \caption{Пластинка чувствительного оттенка.}
        \end{wrapfigure}
        \par Пластинка имеет форму стрелы (рис. 4), вдоль оси которой расположено главное направление, соответствующее большей скорости распространения.
        \par Если пластинка чувствительного оттенка помещена между скрещенными поляроидами и главные направления пластинки не параллельны направлениям разрешенных
        колебаний поляроидов, то при освещении белым светом пластинка кажется окрашенной в лилово-красный цвет. Это объясняется тем, что зеленая компонента линейно
        поляризованного света при прохождении пластинки не меняет поляризации и задерживается вторым поляроидом. Для красной и фиолетовой компонент пластинка создает
        сдвиг фаз, несколько отличный от $2\pi$. На выходе из пластинки красная и фиолетовая компоненты оказываются поэтому эллиптически поляризованными и частично
        проходят через второй поляроид. Таким образом, в известном смысле наблюдаемый в указанном опыте цвет пластинки дополнителен к зеленому.

        \begin{wrapfigure}{l}{0.3\textwidth}
            \centering
            \includegraphics[width=0.3\textwidth]{images/pic5.pdf}
            \caption{К объяснению интерференции поляризованных лучей.}
        \end{wrapfigure}
        \par Если между скрещенными поляроидами поместить пластинку чувствительного оттенка $(\lambda)$ и пластинку в $\lambda/4$ так, чтобы их главные направления
        совпадали, цвет пластинки изменится. Если у пластинки чувствительного оттенка и пластинки в $\lambda/4$ совпадут главные направления, соответствующие большей
        скорости распространения, то разность хода между $E_{\xi}$ и $E_{\eta}$ для зеленого света составит уже $5\lambda/4$. Это соответствует разности хода в
        $\lambda$ для света с большей длиной волны, т. е. для «более красного» света. При освещении этих пластинок (напомним, что они расположены между скрещенными
        поляроидами) белым светом теперь погасится не зеленая, а красная часть спектра, и проходящий свет будет казаться зеленовато-голубым. Если же главные
        направления, соответствующие большей скорости распространения, у пластинки чувствительного оттенка и у пластинки в $\lambda/4$ окажутся перпендикулярными,
        то проходящий свет приобретет оранжево-желтую окраску (погасится фиолетово-голубая часть спектра).
        \par Изменение цвета позволяет, таким образом, определить, какое из главных направлений пластинки в $\lambda/4$ соответствует большей скорости
        распространения.
        \par\textbf{Интерференция поляризованных лучей.} Тонкие двоякопреломляющие пластинки, помещенные между поляроидами, кажутся окрашенными. Эта окраска может
        быть истолкована как результат интерференции поляризованных лучей. На рис. 5 представлена схема для случая скрещенных поляроидов.
        \par Здесь $p_1{p'}_1$ --- разрешенное направление колебаний поляризатора (первого поляроида); $\xi$, $\eta$ --— координатная система, связанная с главными
        направлениями двоякопреломляющей пластинки; $p_2{p'}_2$ --— разрешенное направление колебаний анализатора (второго поляроида). Волны $E_{\xi}$ и $E_{\eta}$
        на выходе из пластинки когерентны, но не могут интерферировать, так как $E_{\xi}\perp E_{\eta}$ . Волны $E_1$ и $E_2$ на выходе второго поляроида также
        являются когерентными и к тому же поляризованы в одной плоскости. Эти волны интерферируют между собой. Результат интерференции определяется зависящим от
        длины волны сдвигом фаз между $E_1$ и $E_2$. В результате интерференции поляризованных лучей пластинка, освещаемая белым светом, кажется окрашенной.
        \par Если поворачивать двоякопреломляющую пластинку, расположенную между скрещенными поляроидами, то соотношение амплитуд волн $E_1$ и $E_2$ и разность фаз
        между ними не изменяются. Это означает, что цвет пластинки при ее поворотах не меняется, а меняется только интенсивность света. За один оборот пластинки
        интенсивность четыре раза обращается в нуль, —-- это происходит при совпадении главных направлений $\xi$ и $\eta$ с разрешенными направлениями колебаний
        поляроидов.
        \par Если же двоякопреломляющую пластинку оставить неподвижной, а второй поляроид повернуть так, чтобы разрешенные направления $p_1{p'}_1$ и $p_2{p'}_2$
        совпали, то волны $E_1$ и $E_2$ приобретают дополнительный фазовый сдвиг на $\pi$ для всех спектральных компонент; поэтому цвет пластинки изменится на
        дополнительный.

    \section{Результаты измерений и обработка данных}

        \subsection{Определение разрешённых направлений поляроидов}

        \par Разместим на оптической скамье осветитель $S$, поляроид $_1$ и чёрное зеркало (пластинку чёрного стекла) так, чтобы плоскость падения была горизонтальна
        (рис. 6). Свет, отражённый от зеркала, будем рассматривать сбоку, расположив глаз таким образом, чтобы вблизи оси вращения зеркала можно было увидеть
        изображение диафрагмы осветителя.

        \begin{wrapfigure}{r}{0.35\textwidth}
            \centering
            \includegraphics[width=0.35\textwidth]{images/pic6.pdf}
            \caption{Определение разрешённого направления поляроида.}
        \end{wrapfigure}
        \par Поворачивая поляроид вокруг направления луча, добейтёмся наименьшей яркости отражённого пятна. Оставим поляроид в этом положении и вращением зеркала
        вокруг вертикальной оси снова добейтёмся минимальной интенсивности отражённого луча. При повороте зеркала необходимо перемещать глаз по горизонтали, чтобы
        следить за отражённым, а не рассеянным пучком. Уточним положения поляроида и зеркала, соответствующие минимуму интенсивности, и определим разрешённое
        направление поляроида.
        \par Разрешённое направление второго поляроида определим, скрестив поляроиды: после поляроида с известной поляризацией поставим второй поляроид и, глядя
        навстречу лучу, вращением второго поляроида добейтёмся минимальной яркости луча. Запишем отсчёт по лимбу поворота поляроидов
        $\phi_1=170^{\circ}\pm5^{\circ}$, $\phi_2=50^{\circ}\pm5^{\circ}$.

        \subsection{Определение угла Брюстера для эбонита}

        \par Поставим на скамью вместо чёрного зеркала (рис. 6) эбонитовую пластину с круговой шкалой.
        \par Повернём эбонитовое зеркало вокруг вертикальной оси так, чтобы его плоскость была перпендикулярна лучу, и попытаемся совместить отражённое от эбонита
        пятно с отверстием осветителя (это можно сделать только в полной темноте, поэтому установим перпендикулярность на глаз). Отметим начало отсчёта по лимбу
        $\phi_0=87^{\circ}$.
        \par Установим направление разрешённых колебаний поляроида $P_1$ горизонтально и найдём угол поворота эбонита $\phi_1=146^{\circ}\pm1^{\circ}$, при
        котором интенсивность отражённого луча минимальна. Значит угол Брюстера равен $\phi_{\text{Б}}=59^{\circ}\pm1^{\circ}$.
        \par Рассчитаем по формуле (3) показатель преломления эбонита. Получим $n=1.66 \pm 0.07$.
        \par Повторим измерения, добавив светофильтр Ф. Измерения дали следущие значения $\phi_0=87^{\circ}$, $\phi_1=145^{\circ}\pm1^{\circ}$. Значит
        $\phi_{\text{Б}}=58^{\circ}\pm1^{\circ}$, и по формуле (3) $n=1.60 \pm 0.06$. Значения совпали в пределах погрешности.
        \par Согласно справочной таблице показатель преломления Эбонита назодится в пределах 1.6-1.7, что согласуется с нашими результатами.

        \subsection{Исследование стопы}

        \begin{wrapfigure}{r}{0.35\textwidth}
            \centering
            \includegraphics[width=0.35\textwidth]{images/pic7.pdf}
            \caption{Исследование стопы.}
        \end{wrapfigure}
        \par Поставим стопу стеклянных пластинок вместо эбонитового зеркала (рис. 7) и подобрём для неё такое положение, при котором свет падает на стопу под углом
        Брюстера. Осветим стопу неполяризованным светом и, рассматривая через поляроиды свет, отражённый от стопы, определим ориентацию вектора $\mathbf{E}$ в
        отражённом луче; затем определим характер поляризации света в преломлённом луче: отражённый луч поляризован в вертикальной плоскости, преломлённый
        поляризован нелинейно, так как при повороте поляроида он становится тусклее, но не исчезает полностью.

        \subsection{Определение главных плоскостей двоякопреломляющих пластин}

        \begin{wrapfigure}{r}{0.3\textwidth}
            \centering
            \includegraphics[width=0.3\textwidth]{images/pic8.pdf}
            \caption{Определение главных направлений в пластинках.}
        \end{wrapfigure}
        \par Поставим кристаллическую пластинку между скрещенными поляроидами $P_1$ и $P_2$ (рис. 8).
        \par Вращая пластинку вокруг направления луча и наблюдая за интенсивностью света, проходящего сквозь второй поляроид, определим, при каком угле по лимбу
        главные направления пластинки совпадают с разрешёнными направлениями поляроидов. Получим $\phi_1=56\pm1$.
        \par Повторяя опыт для второй пластинки получим угол равный $\phi_2=74 \pm 1$
        
        \subsection{Выделение пластин $\lambda/2$ и $\lambda/4$}

        \par Добавим к схеме, изображённой на рис. 8, зелёный фильтр; установим разрешённое направление поляроида горизонтально, а главные направления исследуемой
        пластинки --- под углом $45^{\circ}$ к горизонтали.
        \par С помощью второго поляроида установите, какую поляризацию имеет свет, прошедший пластинку: круговую или линейную с переходом в другой квадрант.
        \par Для первой пластинки освещённость менялась, но не исчезала полностью, значит поляризация была не линейной, а круговой и толщина пластинки равна
        $\lambda/4$.
        \par Для второй пластинки при поо=воротах поляроида свет пропадал полностью, значит поляризация была линейной, и толщина пластинки равна $\lambda/2$.

        \subsection{Определение направлений большей и меньшей скорости в пластинке $\lambda/4$}

        \par Поставим между скрещенными поляроидами пластинку чувствительного оттенка ($\lambda$ для зелёного света), имеющую вид стрелки. Световой вектор,
        ориентированный вдоль направления стрелки, проходит с большей скоростью, перпендикулярный --- с меньшей.

        \begin{wrapfigure}{r}{0.3\textwidth}
            \centering
            \includegraphics[width=0.3\textwidth]{images/pic9.pdf}
            \caption{Определение направлений большей и меньшей скорости.}
        \end{wrapfigure}
        \par Установим разрешённое направление первого поляроида горизонтально и убедимся с помощью второго поляроида, что эта пластинка не меняет поляризацию
        зелёного света в условиях предыдущего опыта. Это связано с тем, что зелёный свет задерживается вторым поляроидом, а красная и синяя компоненты проходят.
        \par Уберём зелёный фильтр и поставим между скрещенными поляроидами пластинку $\lambda$ (стрелка под углом $45^{\circ}$ к разрешённым направлениям
        поляроидов). Глядя сквозь второй поляроид на стрелку, убедимся, что она имеет пурпурный цвет (зелёный свет задерживается вторым поляроидом, а красная и синия
        компоненты проходят).
        \par Добавьте к схеме пластинку $\lambda/4$ (рис. 9), главные направления которой совпадают с главными направлениями пластины $\lambda$ и ориентированы под
        углом $45^{\circ}$ к разрешённым направлениям скрещенных поляроидов.
        \par При повороте рейтера со стрелкой на $180^{\circ}$ вокруг вертикальной оси цвет стрелки меняется от зелёно-голубого до оранжево-жёлтого. При совпадении
        быстрой оси стрелки и главного направления пластинки суммарная разность хода для зеленого света составит $5\lambda/4$, что приблизительно равно длине волны
        в красном спектре, поэтому красный свет будет задерживаться поляроидами, что объясняет зелёно-голубой цвет стрелки. При совпадении медленой оси с главным
        направлением пластинки разность хода $3\lambda /4$ будет составлять длину волны для фиолетово-голубой части спектра, поэтому в результате увидим
        оранжево-жёлтый цвет.

        \subsection{Интерференция поляризованных лучей}
        \begin{wrapfigure}{r}{0.1\textwidth}
            \centering
            \includegraphics[width=0.1\textwidth]{plots/plot1.pdf}
            \includegraphics[width=0.1\textwidth]{plots/plot2.pdf}
            \caption{}
        \end{wrapfigure}
        
		\par Расположим между скрещенными поляроидами мозаичную слюдяную пластинку. Она собрана из 4-х узких полосок слюды, лежащих по сторонам квадрата (две
        полоски «толщиной» $\lambda/4$ и по одной --- $\lambda/2$ и $3\lambda/4$).
        \par В центральном квадратике слюды нет. Главные направления всех пластинок ориентированы параллельно сторонам квадрата.
        \par Вращая пластинку, заметим, что интенсивность света меняется с периодом $\pi/4$.
        \par Теперь, не трогая пластинки, будем вращать второй поляроид. Получим следующие два варианта изображения (рис. 10), повторяющиеся с периодом $\pi/2$.

    \section{Обсуждение результатов}

        \par Большая часть опытов в данной работе носила качественный характе, поэтому нельзя сказать о качетве их проведения. Однако измерение показателя
        преломеления эбонита получилось выполнить с неплохой точностью порядка 5\%.

    \section{Вывод}

        \par Было исследованно много свойств связанных с поляризацией света.

    \newpage\pagenumbering{gobble}

    \begin{figure}[H]
        \begin{center}
            \includegraphics[width=1.4\linewidth, angle=90]{data/approved.pdf}
        \end{center}
    \end{figure}

\end{document}
